\documentclass[12pt]{article}
\usepackage[utf8]{inputenc}
\usepackage[T1]{fontenc}
\usepackage{hyperref}
\title{Tasmanian, North Midlands Language Resources for K-12 Curriculum Development}
\author{Compiled by Claude Sonnet 5.5 (claude-sonnet-5-5) for Andr\'as Kornai}
\date{}
\begin{document}
\maketitle

\section{Introduction}
North Midlands Tasmanian (ISO 639-3 xph, Glottocode neta1235) is an extinct language of Tasmania. The search results show it is attested only in word lists, with a divergent Tamar River variety sometimes called Tyerrernotepanner. The searches found no primer, curriculum or revival programme specific to it.

\section{Primary Educational Resources}
None found. The only living Tasmanian language is the reconstructed palawa kani (\url{https://tacinc.com.au/programs/palawa-kani/}), which draws on Northeastern and Eastern sources. The Tasmanian Aboriginal Centre controls access to its learning materials.

\section{Linguistic and Reference Materials}
\begin{itemize}
\item Tasmanian languages overview: \url{https://en.wikipedia.org/wiki/Tasmanian_languages}.
\item Northeastern Tasmanian languages: \url{https://en.wikipedia.org/wiki/Northeastern_Tasmanian_languages}.
\item Plomley (1976) word-list compilation, catalogued at \url{https://glottolog.org/resource/reference/id/20405}.
\item Joseph Milligan's word lists are among the most extensive. I could not confirm from the results that they cover the North Midlands specifically.
\item A Queen Victoria Museum PDF on Tasmanian tribes: \url{https://www.qvmag.tas.gov.au/files/assets/qvmag/library/publications/occasional/tas-tribes.pdf}.
\item A University of Tasmania eprint found in the search: \url{https://eprints.utas.edu.au/19369/}. Its content was not verified.
\end{itemize}

\section{Religious and Community Resources}
The Tasmanian Aboriginal Centre's palawa kani programme is the community body: \url{https://tacinc.com.au/programs/palawa-kani/}. No religious texts were found.

\section{Digital and Audio Resources}
None found.

\section{Curriculum Development Context}
No usable teaching material exists for this language itself. A heritage unit would be about Tasmanian Aboriginal history and languages, in partnership with the Tasmanian Aboriginal community.

\end{document}